\problema{Top K Frequent Elements}{347}{src/04-heaps-topk/347-top-k-frequent-elements.cpp}{M}

Nos dan un arreglo \texttt{nums} y un entero \texttt{k}. Hay que devolver los
\texttt{k} valores que aparecen con mayor frecuencia en \texttt{nums}. El
orden de la respuesta no importa, y se garantiza que la respuesta es única
(no hay empates ambiguos justo en la posición \texttt{k}).

\paragraph{La idea, con un ejemplo de la vida real.} Piensa en la lista de
canciones más reproducidas de un usuario en un servicio de música: no
importa \emph{cuándo} se reprodujo cada canción, solo cuántas veces se
reprodujo cada una. Para armar el top \texttt{k}, primero cuentas cuántas
veces se repite cada canción, y luego agrupas las canciones en
``casilleros'' según su número de reproducciones: un casillero para las que
se reprodujeron una vez, otro para las que se reprodujeron dos veces, y así
sucesivamente. Como el número máximo de reproducciones posible es el total
de reproducciones registradas, hay un número limitado y conocido de
casilleros. Para armar el top \texttt{k} basta recorrer los casilleros
empezando por el de más reproducciones y bajando, tomando canciones hasta
completar \texttt{k}.

\begin{center}
\ttfamily
\begin{tabular}{l}
nums = [1,1,1,2,2,3], k = 2\\
frecuencias: 1 -> 3, 2 -> 2, 3 -> 1\\
buckets[3]=\{1\}, buckets[2]=\{2\}, buckets[1]=\{3\}\\
recorriendo de mayor a menor: 1, luego 2 -> respuesta \{1, 2\}\\
\end{tabular}
\end{center}

\paragraph{Cómo lo resuelve el código, paso a paso.}
\begin{enumerate}
  \item Contamos cuántas veces aparece cada valor con un mapa de frecuencias.
  \item Creamos un arreglo de \emph{buckets} (casilleros) de tamaño
        $n+1$, donde \texttt{buckets[f]} guardará todos los valores que
        aparecen exactamente \texttt{f} veces.
  \item Recorremos el mapa de frecuencias y colocamos cada valor en el
        casillero que le corresponde según su frecuencia.
  \item Recorremos los casilleros de mayor a menor frecuencia, agregando sus
        valores a la respuesta, hasta juntar \texttt{k} elementos.
\end{enumerate}
Esto se conoce como \emph{bucket sort} (ordenamiento por casilleros): en vez
de comparar elementos entre sí para ordenarlos, usamos directamente el valor
que queremos ordenar (aquí, la frecuencia) como índice de un arreglo. Como
cada valor cae en exactamente un casillero, llenar los casilleros cuesta
$O(n)$, y recorrerlos también.

\lstinputlisting{src/04-heaps-topk/347-top-k-frequent-elements.cpp}

\complejidad{$O(n)$}{$O(n)$}

\paragraph{¿Qué significa esa complejidad?} Tanto contar frecuencias como
llenar y recorrer los casilleros toma un tiempo proporcional al tamaño del
arreglo original, sin ningún logaritmo de por medio: es más rápido que
cualquier solución basada en ordenar ($O(n \log n)$) o incluso que un
min-heap de tamaño \texttt{k} ($O(n \log k)$), aunque ambas alternativas
también son aceptables en una entrevista.

\paragraph{Consejos para la entrevista (Amazon).} Si te preguntan la primera
solución que se te ocurre, es normal mencionar primero el min-heap de
tamaño \texttt{k} (el mismo patrón del problema 215): es más fácil de
explicar y sigue siendo eficiente. El \emph{bucket sort} es la mejora
``inteligente'' que muestra que reconoces cuando el rango de valores
posibles (aquí, la frecuencia máxima es $n$) es lo bastante acotado como
para usarlo directamente de índice. Menciona explícitamente por qué el
tamaño del arreglo de buckets es $n+1$ y no algo arbitrario: la frecuencia
de cualquier valor nunca puede superar el tamaño del arreglo de entrada.
